\documentclass[11pt]{article}
\usepackage[margin=1in]{geometry}
\usepackage{amsmath,amssymb}
\usepackage{graphicx}
\usepackage{booktabs}
\usepackage{hyperref}
\usepackage{microtype}

\title{\textbf{Paper 40: Real Continual Learning in VCML ---\\
Zone Boundary Shift, Forgetting, and Two-Task Interleaving}}
\author{VCSM Theoretical Series}
\date{2026-03-10}

\begin{document}
\maketitle

\begin{abstract}
Paper~38 tested continual learning by inverting wave polarity (suppression $\leftrightarrow$
excitation) mid-run, but this perturbation is invisible to sg4n (polarity-blind by
construction).  Here we test \emph{genuine} continual learning by shifting the spatial
positions of zone boundaries at $T_\mathrm{shift} = 1500$ steps.
Three experiments probe adaptation, forgetting, and multi-task stability:
(A)~zone boundary shift with overlap fractions of 0\% (cyclic permutation, $\Delta=10$)
and 50\% (partial overlap, $\Delta=5$), versus a no-shift reference;
(B)~two-task interleaving from $t=0$, alternating every 500 steps;
(C)~the same interleaving begun after a 2000-step commitment phase for task~A.
Key findings:
(1) a cyclic-permutation shift ($\Delta = z_w$) is \emph{immediately transparent}:
sg4n under the new zone labelling equals sg4n under the old labelling at the moment of
shift, because sg4 is invariant under permutation of zone indices;
(2) a half-width shift ($\Delta = z_w/2$) causes genuine disruption --- sg4n under the
new labelling drops and recovers slowly (still rising at $T=4000$);
(3) prior commitment provides \emph{no long-term advantage} in two-task interleaving:
both tasks converge to the same steady-state sg4n ($\approx 0.29$ and $0.27$) regardless
of whether task~A was committed first.
\end{abstract}

\section{Introduction}

Continual learning asks whether a system can acquire new representations without
destroying previously learned ones.
Paper~38 attempted to test this by inverting the suppression/excitation assignment of all
zones at step $T_\mathrm{shift} = 2000$, but this perturbation is a \emph{polarity flip}
--- it changes the sign of the wave activation without moving zone boundaries.
Because sg4n measures mean pairwise distances between zone-mean fieldM vectors, it is
invariant to a global sign inversion of all zone vectors (the inter-zone distances are
unchanged).
The polarity-flip experiment confirmed that sg4n is polarity-blind, but it could not
detect whether VCML's spatial structure updated.

The correct test is a \emph{zone boundary shift}: move the spatial positions that define
each zone, so that a cell previously in zone~$z$ now belongs to zone~$z'$.
Unlike a polarity flip, this is not invisible to sg4n under the new labelling.
If the system fails to adapt, sg4n under the new labelling stays near zero.
If the system adapts, sg4n under the new labelling recovers.

We study three cases:
\begin{enumerate}
  \item \textbf{Cyclic permutation} ($\Delta = z_w = 10$): all zone boundaries shift by
        one full zone width.  Zone~0 cells become zone~1 cells, zone~1 becomes zone~2,
        etc.\ (modular wrap).  This is a consistent relabelling: each cell goes from
        one zone to exactly one other zone, with no mixing.
  \item \textbf{Half-width shift} ($\Delta = z_w/2 = 5$): zone boundaries shift by half
        a zone width.  Half of every zone's cells cross a boundary into a neighbouring
        zone; the other half remain.  New zones are composed of cells from two different
        old zones.
  \item \textbf{No shift} (reference): standard run, no boundary change.
\end{enumerate}

We also test two-task interleaving: alternating between zone configuration~A
($\Delta = 0$) and configuration~B ($\Delta = 5$) every 500 steps,
comparing the case where interleaving starts immediately ($t=0$) versus after a
2000-step commitment phase.

\section{Methods}

\subsection{Base System}

Standard VCML: $H=40$, $\mathrm{HALF}=40$, $N_\mathrm{zones}=4$, $z_w=10$,
$\mathrm{WR}=4.8$, $\mathrm{FA}=0.16$, $T_\mathrm{shift}=1500$, 5 seeds per condition.
The commitment epoch is $t^* \approx 2000$ (Paper~27), so the shift occurs \emph{before}
full commitment.

\subsection{ShiftableEnv}

A \texttt{ShiftableEnv} environment launches waves at positions consistent with a
zone-boundary offset $\Delta$.
Zone class $z$ in shifted coordinates corresponds to positions $c_x$ in the active region
satisfying $(c_x - \mathrm{HALF} + \Delta) \bmod \mathrm{HALF} \in [z \cdot z_w,
(z+1) \cdot z_w)$.
The attribute \texttt{env.delta} is changed at $T_\mathrm{shift}$;
thereafter all wave launches respect the new boundary positions.

\subsection{Dual-Labelling Metric}

At each checkpoint we compute sg4n under \emph{two} zone labellings simultaneously:
\begin{itemize}
  \item \textbf{sg4n\_new}: using the current ($\Delta$) zone labelling.
        Measures adaptation to the new configuration.
  \item \textbf{sg4n\_old}: always using $\Delta=0$ (original boundaries).
        Measures how much of the old structure survives.
\end{itemize}
Both can be high simultaneously: the system retains old structure while building new.
Both near zero indicates failure to maintain any zone differentiation.

\section{Results}

\subsection{Experiment A: Zone Boundary Shift}

\begin{table}[h]
  \centering
  \caption{sg4n under new labelling at selected time points.
           Shift occurs at $T_\mathrm{shift}=1500$; first post-shift checkpoint is $t=1600$.}
  \label{tab:expa}
  \small
  \begin{tabular}{lccccc}
    \toprule
    Condition & $t=1400$ & $t=1600$ & $t=2000$ & $t=3000$ & $t=4000$ \\
              & (pre-shift) & (post) & & & \\
    \midrule
    C\_ref  ($\Delta=0$)  & 0.326 & 0.375 & 0.243 & 0.259 & 0.252 \\
    C\_half ($\Delta=5$)  & 0.326 & 0.162 & 0.182 & 0.284 & \textbf{0.798} \\
    C\_full ($\Delta=10$) & 0.326 & \textbf{0.415} & 0.236 & 0.161 & 0.194 \\
    \bottomrule
  \end{tabular}
\end{table}

Table~\ref{tab:expa} shows markedly different responses across the three conditions.

\paragraph{Cyclic permutation ($\Delta=10$) is immediately transparent.}
At $t=1600$ (100 steps after the shift), sg4n\_new ($=0.415$) equals sg4n\_old ($=0.415$)
--- they are identical because sg4 is a mean of \emph{pairwise} distances between zone
means, and a cyclic permutation of zone indices only permutes the distance matrix rows and
columns without changing any individual distance.
The system is already in the correct configuration for the new labelling without any
adaptation required.
This extends the alignment-blindness result of Paper~39 to the single-region continual
learning setting: a consistent relabelling (here a permutation rather than a spatial
offset) does not constitute a genuine learning challenge.

After the shift, sg4n for C\_full decays along the normal metastability trajectory
(peak at $t\approx1600$, half-life $\sim800$ steps), suggesting the env change
(waves now drive the shifted configuration) gradually refocuses structure.

\paragraph{Half-width shift ($\Delta=5$) causes genuine disruption and slow adaptation.}
At $t=1600$, sg4n\_new drops to $0.162$ from the pre-shift value of $0.326$ --- a 50\%
reduction.
The disruption arises because new zones are composites of two old zones:
new zone~0 = \{cells 0--4 (old zone~0)\} $\cup$ \{cells 35--39 (old zone~3)\}.
These cells carry fieldM from two different attractor states (suppression- and
excitation-type).
Under the new labelling, within-zone variance rises and between-zone distance falls,
lowering sg4n.

Recovery is slow: sg4n\_new is still rising at $t=4000$ ($=0.798$), far exceeding the
reference trajectory peak ($0.465$).
The eventual high value reflects the system successfully learning the new zone
configuration: waves drive cells at mixed-heritage positions to adopt a unified new zone
identity, and the copy-forward loop propagates this across the zone through births.
The long recovery time ($>2500$ steps) is consistent with the commitment-epoch timescale:
re-differentiation under the new zone map requires a new commitment-epoch cycle starting
from the disrupted state.

\paragraph{Old structure decays on the metastability timescale.}
For all shifted conditions, sg4n\_old evolves similarly to the reference (Figure 1b):
it does not collapse immediately at the shift but decays over $\sim800$--$1600$ steps
(the half-life from Paper~38).
The copy-forward loop preserves old zone structure even while new waves drive toward the
new configuration.
The old and new structures coexist transiently.

\subsection{Experiment B: Interleaved Switching Without Prior Commitment}

Alternating between task~A ($\Delta=0$) and task~B ($\Delta=5$) every 500 steps from $t=0$
produces a systematic asymmetry:
task~A converges to a higher steady-state sg4n ($\approx 0.48$) than task~B ($\approx
0.30$), a ratio of 1.61.

The asymmetry arises from task ordering: task~A is always the first task in each 500-step
period.
The first 500 steps (task~A only) establishes structure from a cold start.
When task~B first appears at $t=500$, it must build structure against the already-formed
task~A fieldM, which the copy-forward loop partially perpetuates across each switch.
The structural head-start of task~A is never fully overcome.

\subsection{Experiment C: Post-Commitment Interleaved Switching}

Task~A is run alone for 2000 steps (commitment phase, sg4n\_A peak $=0.453$), then
interleaving begins.

The result is convergence: by $T=5000$, sg4n\_A $= 0.294$ and sg4n\_B $= 0.269$,
a ratio of $1.09$.
Prior commitment to task~A provides only a transient advantage.
The committed attractor for task~A is eroded by the periodic task~B episodes
(erosion of $\approx35\%$ over the interleaving period), while task~B builds from zero
to nearly the same level.
The system settles into a steady state where both tasks coexist at roughly equal strength.

This convergence implies that the VCML attractor is not an absorbing state.
The committed structure from $t^*\approx2000$ provides resistance proportional to the
metastability half-life ($\sim800$ steps) but is eventually overwritten by the
interleaved input.
The system does not exhibit catastrophic forgetting of task~A, but it does exhibit
\emph{gradual forgetting} at a rate determined by the turnover and field decay dynamics.

\section{Discussion}

\subsection{The Three Modes of Zone-Shift Response}

Zone boundary shifts reveal three qualitatively distinct behaviours:

\textbf{Consistent permutation} ($\Delta = k \cdot z_w$, integer $k$): no learning
challenge.
sg4 is permutation-invariant; the old structure is immediately readable under the new
labelling.
This is not a continual learning scenario but a relabelling scenario.

\textbf{Partial overlap} ($0 < \Delta < z_w$): genuine learning challenge.
The old and new configurations are incompatible: some cells must switch zone identity.
Recovery requires a full re-commitment cycle ($>2000$ steps).
The copy-forward loop contributes to both the disruption (propagating old fieldM across
new zone boundaries) and the recovery (propagating new zone identity through births).

\textbf{No shift}: the metastability trajectory (Paper~38) is the baseline.
Structure peaks near $t^*$ and decays with half-life $\sim800$ steps.

\subsection{Committed Structure Is Not Rigid}

Paper~38 established a bimodal attractor with $P_\mathrm{high}(\mathrm{geo}) = 0.50$ and
$G_\mathrm{basin} = 1.88$.
The current experiments show that once in the high-structure basin, the system is not
permanently trapped there.
Two-task interleaving with a 500-step period drives the system out of the task~A basin
and toward a mixed steady state, at a rate determined by the period length relative to the
metastability half-life.

This has direct implications for biological plasticity.
In the hippocampal dentate gyrus, mandatory neurogenesis (turnover) prevents crystallisation
of zone boundaries, maintaining adaptability at the cost of gradual forgetting.
The 500-step interleaving period here corresponds to a biological learning schedule
where the input environment alternates faster than the metastability half-life.
When the switching period exceeds the half-life, old structure survives; when it falls
below, the system tracks the current input.

\subsection{Task-Ordering Asymmetry and the First-Mover Advantage}

Experiment~B reveals a first-mover advantage: the task presented first in each interleaving
cycle maintains higher sg4n.
This asymmetry is structural, not parametric --- both tasks use identical wave parameters
and the same FA.
The advantage arises because the copy-forward loop propagates task~A fieldM during the
500-step task~A phase, which the 500-step task~B phase must partially overwrite.
If the periods were equal and the system were linear, the advantage would cancel.
The asymmetry implies that the copy-forward loop introduces a non-linearity:
established fieldM resists overwriting more than empty fieldM resists writing.

\section{Conclusion}

Three conclusions close the continual learning thread opened in Paper~38:

\begin{enumerate}
  \item \textbf{Permutation-invariance prevents detection of relabelling.}
        A cyclic zone shift ($\Delta = z_w$) is immediately transparent.
        What Paper~38 called ``polarity blindness'' generalises:
        any \emph{consistent relabelling} of zones leaves sg4 unchanged.
        The genuine learning challenge requires \emph{spatial reassignment with overlap}.

  \item \textbf{Partial overlap causes disruption with slow recovery.}
        A half-width shift ($\Delta = z_w/2$) drops sg4n\_new by 50\% at the shift,
        then requires $>2500$ steps to recover --- comparable to the initial
        commitment timescale.
        Old structure and new structure coexist during this recovery period.

  \item \textbf{Prior commitment provides no permanent advantage under interleaving.}
        Two tasks converge to equal steady-state sg4n when the interleaving period
        ($500$ steps) is shorter than the commitment epoch ($\sim2000$ steps).
        The system exhibits gradual, not catastrophic, forgetting;
        the decay rate is set by the metastability half-life ($\sim800$ steps).
\end{enumerate}

\end{document}
